% Reconstructed after workspace maintenance from the retained mathematical
% derivation.  Original pre-maintenance SHA256:
% f354724c43b98d2fc54119d7fb74604b4e47306919e348528fb6f35924f8cb3a.
% Byte identity is not asserted; current source is independently rechecked.
\section{The central physical clock and exact section observations}
\label{sec:cp-central}
The geometric boundary-layer result gives actual central comparison
records outside a specified positive original source.  The remaining
clock comparison must retain its two endpoint work terms.  Dropping
them would turn a variational identity into a false estimate for the
physical roof.

Keep the notation and original family of Section~\ref{sec:bl-global}.
Fix integers $d\ge1$, $2d<m$, and put $b=m-d$.  In this section
$Y$ is a global disk-word minimum whose consecutive vertices
$Y_d,\ldots,Y_b$ form an actual regular physical subrecord, with the
same consecutive centers as $Q_d,\ldots,Q_b$.  The complement of the
source in Theorem~\ref{thm:tb-uniform-source}, for an applicable
$D\le d$, supplies this physicality.  The statements below need only
the stated physical subrecord, not a physical continuation to $Y$ or
regularity of either discarded end portion.  These are comparisons
inside the original record; they do not replace the original full
conditional event.

Write
\begin{equation}\label{eq:cp-central-objects}
 \begin{gathered}
 I_d=\{d,\ldots,b\},\qquad J_d=\{d,\ldots,b-1\},\qquad m_d=m-2d,\\
 L_d(Q)=\sum_{i\in J_d}\tau_i,\quad
 L_d(Y)=\sum_{i\in J_d}\ell_i,\quad
 W_d=v_{b-1}\cdot\delta_b-v_d\cdot\delta_d,\\
 \mathsf T_{{\rm clock},d}=L_d(Y)-L_d(Q)-W_d .
 \end{gathered}
\end{equation}

\begin{theorem}[The exact central clock with its boundary work]
\label{thm:cp-clock}
For every such original record and minimum,
\begin{equation}\label{eq:cp-clock}
 \mathsf T_{{\rm clock},d}
   =\frac1R\sum_{i=d+1}^{b-1}H_i
                     +\frac12\sum_{i=d}^{b-1}F_i\ge0.
\end{equation}
In particular, with the same $K_{\rm tail}$ as in
Theorem~\ref{thm:bl-tail},
\begin{equation}\label{eq:cp-clock-cost}
 0\le\mathsf T_{{\rm clock},d}\le K_{\rm clock}d^{-3},
 \qquad
 K_{\rm clock}=
 \left(\frac{32r}{R_-}+64\right)
                    \frac{\sqrt{K_{\rm tail}}}{\tau_-}.
\end{equation}
The length on either side is the unit-speed physical time of that
central subrecord.  No clock is averaged or rescaled.
\end{theorem}
\begin{proof}
For each included flight, the identity between two unit directions is
\[
 \ell_i-\tau_i
  =v_i\cdot(\delta_{i+1}-\delta_i)
                                +\tfrac12\ell_i|e_i|^2 .
\]
Sum it over $d\le i<b$.  Finite integration by parts gives boundary
term $W_d$ and interior coefficients $v_{i-1}-v_i=-2c_in_i$.
Both actual contact points $Q_i,Y_i$ lie on the same circle, so
\[
 |Q_i-O_i+\delta_i|^2=R^2
 \quad\Longrightarrow\quad
 2Rn_i\cdot\delta_i+|\delta_i|^2=0 .
\]
Thus each interior coefficient contributes $H_i/R$ with a positive
sign.  This proves the exact identity, including the case of one
central flight, when its contact sum is empty.  Apply the contact
and flight estimates of Corollary~\ref{cor:bl-energy}, and enlarge
$R^{-1}$ to $R_-^{-1}$, to obtain the displayed constant.  This
argument uses circular equality only at the included physical
contacts.  The disk support inequality for a nonphysical interior
vertex would not justify the same clock equality.
\end{proof}

Let $x_i=(\alpha_i,p_i)$ and $\widehat x_i$ be the original and
comparison collision states on $I_d$.  Use outgoing tangential
velocity at $i<b$ and the incoming tangential projection at $b$.
Specular reflection makes that terminal projection the actual
collision coordinate without requiring a next flight.  The distance
$d_\infty$ is the maximum of the shortest circular angle distance
and the momentum distance.  Define
\begin{equation}\label{eq:cp-phase-cost}
 \mathsf T_{{\rm phase},d}
      =\sum_{i=d}^{b}d_\infty(x_i,\widehat x_i)^4,\qquad
 K_{\rm phase}
      =\frac{8192\sqrt{K_{\rm tail}}}{\tau_-^2}
                                      +\frac{24K_{\rm tail}}{R_-^4}.
\end{equation}

\begin{theorem}[A complete central phase and section payment]
\label{thm:cp-phase}
One has
\begin{equation}\label{eq:cp-phase-bound}
 \mathsf T_{{\rm phase},d}\le K_{\rm phase}d^{-3}.
\end{equation}
Let $\eta_R=\one_{Y_R^*}$ be exactly the original section indicator.
For $\rho>0$ put
\[
 B_{\rho,d}(Q)=\sum_{i=d}^b
       \one_{\{\operatorname{dist}_\infty(x_i,\partial Y_R^*)\le\rho\}},
 \qquad
 N_{{\rm cut},d}(Q,Y)=\sum_{i=d}^b|\eta_R(x_i)-\eta_R(\widehat x_i)|.
\]
Then, for the actual observations on both subrecords,
\begin{equation}\label{eq:cp-section-bound}
 N_{{\rm cut},d}
 \le B_{\rho,d}+\rho^{-4}\mathsf T_{{\rm phase},d}
 \le B_{\rho,d}+\rho^{-4}K_{\rm phase}d^{-3}.
\end{equation}
Occupation uses only $J_d$; the last section test at $b$ is separate.
In particular the absolute occupation difference is bounded by the
sum in \eqref{eq:cp-section-bound} with its last summand removed.
\end{theorem}
\begin{proof}
The circle estimate in Theorem~\ref{thm:ss-section} applies contact
by contact:
\[
 d_\infty(x_i,\widehat x_i)^4
 \le 8|e_{\iota(i)}|^4+24R_-^{-4}|\delta_i|^4,\qquad
 \iota(i)=
 \begin{cases}i,&i<b,\\ b-1,&i=b.\end{cases}
\]
Indeed angle distance is at most $2|\delta_i|/R_-$ and tangential
momentum distance at most $|e_{\iota(i)}|+|\delta_i|/R_-$.
The last edge occurs twice, including when the subrecord has only
one edge.  All other included edges occur once.  Since $|e_i|\le2$
and $\ell_i\ge\tau_-$, the sum of their direction contributions is
at most
\[
 8\sum_{i\in I_d}|e_{\iota(i)}|^4
                  \le\frac{64}{\tau_-}\sum_{i=d}^{b-1}F_i .
\]
Use Corollary~\ref{cor:bl-energy} for this sum and
Theorem~\ref{thm:bl-tail} for the position sum.  Enlarging
$(d+1)^{-3}$ to $d^{-3}$ gives exactly $K_{\rm phase}$.

A shortest product-coordinate path between states with different
section indicators meets the original section boundary.  Therefore
a discrepancy is either in its closed $\rho$ collar or has
$d_\infty(x_i,\widehat x_i)>\rho$.  Its indicator in the latter
case is at most $\rho^{-4}d_\infty(x_i,\widehat x_i)^4$.
Sum over every included contact.  Boundary states have been counted
explicitly; no null-set deletion of their weighted source is used.
Dropping the last nonnegative summand gives the occupation claim.
\end{proof}

\begin{proposition}[All central Fourier marks with the exact clock phase]
\label{prop:cp-twists}
The central collision count $m_d$ and lattice displacement
$k_d=O_b-O_d$, expressed in the original lattice basis, are identical
on the two physical subrecords.  Put
$A_d(Q)=\sum_{i=d}^{b-1}\eta_R(x_i)$ and define $A_d(Y)$ in the same
way.  For every $\xi\in\mathbb R^2$ and $s,\zeta\in\mathbb R$,
\begin{align}\label{eq:cp-exact-twist}
 &\left|e^{i[\xi\cdot k_d+sL_d(Y)+\zeta A_d(Y)]}
       -e^{isW_d}e^{i[\xi\cdot k_d+sL_d(Q)+\zeta A_d(Q)]}\right|
       \notag\\
 &\hspace{12mm}\le
 |s|\mathsf T_{{\rm clock},d}
                  +2|\sin(\zeta/2)|\,N_{{\rm cut},d}.
\end{align}
If the two central endpoint section tests are inserted, their
additional error is at most the sum of those two actual indicator
discrepancies.  All these comparisons retain every central collision.
\end{proposition}
\begin{proof}
The ordered centers give $k_d$ and $m_d$ exactly.  Substitute
\eqref{eq:cp-clock} into the left side, retaining the phase
$e^{isW_d}$ of modulus one.  The clock error is bounded by
$|e^{is\mathsf T_{{\rm clock},d}}-1|\le
 |s|\mathsf T_{{\rm clock},d}$.  For an integer $q$, telescoping
$q$ unit increments gives
$|e^{i\zeta q}-1|\le2|q||\sin(\zeta/2)|$; it also holds for
negative $q$.  Take $q=A_d(Y)-A_d(Q)$ and use the preceding
theorem.  The product of two indicators changes by at most the sum
of their two changes, which proves the endpoint assertion.
\end{proof}

\begin{remark}\label{rem:cp-full-event}
The full source event is still $\Pi_{n,k,m,R}$ and its roof still
$L_{m,R}(Q)$.  Neither the central occupation nor the phase-corrected
central length is substituted into that event.  The end portions,
the section collar and the physicality-exception source remain
present in the full problem until their own required estimates have
been proved.  In particular $\mathsf T_{{\rm clock},d}$ alone
cannot bound $|L_d(Y)-L_d(Q)|$: the term $W_d$ is indispensable.
\end{remark}
